Table~\ref{tab:data} summarizes the panel; the full per-model results are
in the supplement. Three patterns are consistent across the four systems.

\paragraph{Calendar gains replicate everywhere (H1).}
The best calendar-augmented model lowers held-out test MAE relative to
the best internal-history model in every city: $-11.9\%$ in New York
(207.0 $\rightarrow$ 182.4), $-14.0\%$ in Chicago (59.2 $\rightarrow$
50.9), $-8.6\%$ in San Francisco (34.5 $\rightarrow$ 31.5), and
$-10.2\%$ in Austin (16.5 $\rightarrow$ 14.8). All four differences are
significant under a paired day-block bootstrap ($p \le 0.0005$;
supplement). Against the same-weekday seasonal-naive baseline, the
selected LightGBM calendar+weather model reduces test MAE by 41.4\% (New
York, 282.4 $\rightarrow$ 165.4), 41.5\% (Chicago), 31.3\% (San
Francisco), and 31.3\% (Austin). H1 survives in all four systems.

\paragraph{Weather's marginal value is real but city-dependent.}
On top of calendar features, target-day weather further reduces MAE by
9.3\% in New York ($p = 0.0005$), 3.3\% in Chicago ($p = 0.0015$), and
4.7\% in San Francisco ($p = 0.0010$) --- and by exactly nothing in
Austin ($+0.001$ MAE, $p = 0.995$). The Austin null is retained as a
finding: a public signal that helps one city is not guaranteed to help
another, even within one country and one data standard. Moreover, the
lagged-weather sensitivity variant (assumption A3 removed) erases most
of New York's weather gain (178.6 versus 165.4 with target-day weather,
against a 182.4 calendar-only reference), so much of the weather value
is conditional on an accurate day-ahead weather forecast --- a realistic
but explicitly stated operational premise.

\paragraph{Model ordering is stable.}
LightGBM is the best point forecaster in every city and feature set,
followed by the random forest; ridge trails both, and both naive
baselines are dominated (their rows remain in all tables). Five
expanding-window rolling-origin folds reproduce the feature-set ordering
(internal $>$ calendar $>$ calendar+weather in MAE) in every city, so
the headline comparisons are not artifacts of a single split.

\paragraph{Pooling across cities mostly hurts (RQ4).}
A single pooled model with city indicators \emph{raises} test MAE
relative to local models in New York ($+1.8\%$, $p = 0.038$), Chicago
($+6.6\%$, $p = 0.001$), and Austin ($+2.4\%$, n.s.), and helps only in
San Francisco ($-2.0\%$, n.s.). Zero-shot transfer (training with the
target city held out entirely) degrades MAE between $+44\%$ (Chicago)
and $+147\%$ (New York), except in San Francisco ($+3.6\%$). With six
years of local history per city, negative transfer is the norm in this
benchmark; the negative result is kept and discussed in Section~9.

\paragraph{Probabilistic forecasts are sharp but under-dispersed.}
The quantile LightGBM's central 90\% intervals cover 76.8--81.2\% of
test outcomes (50\% intervals: 34--41\%) across cities. We report this
miscalibration rather than re-tuning it away post hoc; Section~7 shows
the uncertainty information is operationally valuable despite it, and
conformal recalibration is noted as future work.

\begin{figure}[t]\centering
\includegraphics[width=0.9\linewidth]{fig2_accuracy_gain.pdf}
\caption{Best held-out test-MAE reduction versus the trailing-mean naive
baseline, by feature set and city.}
\label{fig:gain}
\end{figure}
\begin{figure}[t]\centering
\includegraphics[width=0.95\linewidth]{fig4_fold_stability.pdf}
\caption{Rolling-origin fold MAE (LightGBM) by feature set: the
calendar and weather orderings are stable across folds in all cities.}
\label{fig:folds}
\end{figure}
